\section{Implementación de la primera entrega}
\label{sec:implementacion}

\subsection{Alcance}

Se implementaron en C++17 orientado a objetos, con ZeroMQ, los componentes de la
arquitectura base con las operaciones de \textbf{consulta de eventos del mes} y
\textbf{reserva de cupos}:

\begin{table}[H]
\centering
\caption{Componentes implementados.}
\begin{tabularx}{\textwidth}{p{3.8cm}p{4.6cm}Y}
\toprule
Componente (Figura~\ref{fig:arquitectura-basica}) & Ejecutable & Clases \\
\midrule
Clientes & \codigo{bin/cliente} & \codigo{Cliente}, \codigo{ConsolaCliente} \\
Gestor de Transacciones & \codigo{bin/gestor-transacciones} & \codigo{GestorTransacciones} \\
Procesamiento RMC & \codigo{bin/procesador-rmc} & \codigo{ProcesadorRMC} \\
Servicio de Persistencia & \codigo{bin/persistencia} & \codigo{ServicioPersistencia},
\codigo{BaseDatos}, \codigo{AlmacenamientoArchivo}, modelo de dominio \\
Común & — & \codigo{Servicio}, \codigo{Mensaje}, \codigo{ClienteConfiable} \\
\bottomrule
\end{tabularx}
\end{table}

El componente de Cancelación de reservas está diseñado (sección~\ref{sec:diseno})
pero no se implementó, porque esta entrega solo incluye consulta y reserva. Aunque
la tolerancia a fallos no se exigía, el código ya incluye los reintentos
\emph{Lazy Pirate}, los latidos entre gestor y procesadores y las reservas
idempotentes, porque son la base de la segunda entrega.

Persistencia guarda los datos en memoria y escribe cada reserva en un archivo de
bitácora antes de responder; al arrancar vuelve a aplicar la bitácora. Arranca con
datos de ejemplo: cuatro eventos en noviembre y diciembre de 2026 (en el Teatro
Mayor, una sala de 5 cupos, un auditorio y un lugar externo) y los estudiantes
\codigo{EST-1}, \codigo{EST-2} y \codigo{EST-3}.

\subsection{Compilación y ejecución}

\begin{lstlisting}
sudo apt install g++ make libzmq3-dev cppzmq-dev
make
scripts/iniciar_sistema.sh base      # o "escalable" (3 procesadores RMC)
bin/cliente EST-1                    # menu: 1 consultar mes, 2 reservar
scripts/detener_sistema.sh
\end{lstlisting}

Para ejecutar en varias máquinas se indican las direcciones con variables de
entorno (\codigo{GESTOR\_ENDPOINT}, \codigo{GESTOR\_BACKEND},
\codigo{PERSISTENCIA\_ENDPOINT}, etc.). El código fuente y el README están en el
repositorio: \repositorio.

\subsection{Verificación realizada}

\begin{itemize}
    \item Consulta del mes \codigo{2026-11}: lista las fechas con cupos y precios.
    \item Errores: fecha inexistente, cupos en 0, mes mal escrito, estudiante no
          registrado y reserva duplicada.
    \item Tres estudiantes pidieron 2 cupos al mismo tiempo en la sala de 5: dos
          reservas se crearon y la tercera recibió \codigo{CUPOS\_INSUFICIENTES}.
    \item Tras reiniciar persistencia, las reservas se conservaron.
\end{itemize}
